% from ega2-1-fr.tex
\oldpage[II]{102}

Si $f$ est propre, il en est de même de $f\circ j_i$, puisque $j_i$ est une
immersion fermée (\hyperref[II.5.4.2-fr]{5.4.2})~; comme $h_i$ est une
immersion fermée, donc un morphisme séparé, $f_i$ est propre par
(\hyperref[II.5.4.3-fr]{5.4.3}). Supposons inversement que chacun des $f_i$
soit propre, et considérons le préschéma $Z$ somme des $X_i$, soit $u$ le
morphisme $Z\to X$ qui se réduit à $j_i$ dans chacun des $X_i$. La
restriction de $f\circ u$ à chaque $X_i$ est égale à
$f\circ j_i=h_i\circ f_i$, donc est propre puisqu'il en est ainsi de $h_i$
et de $f_i$ (\hyperref[II.5.4.2-fr]{5.4.2})~; il résulte alors aussitôt de
la déf. (\hyperref[II.5.4.1-fr]{5.4.1}) que $f\circ u$ est propre. Mais
comme $u$ est surjectif par hypothèse, on conclut que $f$ est propre par
(\hyperref[II.5.4.3-fr]{5.4.3}).

\begin{corollary}[5.4.6]
\label{II.5.4.6-fr}
Soit $f:X\to Y$ un morphisme séparé et de type fini~; pour que $f$ soit
propre, il faut et il suffit que
$f_{\mathrm{red}}:X_{\mathrm{red}}\to Y_{\mathrm{red}}$ le soit.
\end{corollary}

C'est un cas particulier de (\hyperref[II.5.4.5-fr]{5.4.5}) avec $n=1$,
$X_1=X_{\mathrm{red}}$ et $Y_1=Y_{\mathrm{red}}$
(\hyperref[I.5.1.5-fr]{I, 5.1.5}).

\begin{env}[5.4.7]
\label{II.5.4.7-fr}
Si $X$ et $Y$ sont des préschémas noethériens et $f:X\to Y$ un morphisme
séparé de type fini, on peut, pour vérifier que $f$ est propre, se ramener à
le faire pour des morphismes dominants de préschémas intègres. En effet,
soient $X_i$ ($1\leq i\leq n$) les composantes irréductibles en nombre fini
de $X$ et considérons pour chaque $i$ l'unique sous-préschéma fermé réduit de
$X$ ayant $X_i$ pour espace sous-jacent, que nous désignerons encore par
$X_i$ (\hyperref[I.5.2.1-fr]{I, 5.2.1}). Soit d'autre part $Y_i$ l'unique
sous-préschéma fermé réduit de $Y$ ayant $\overline{f(X_i)}$ pour espace
sous-jacent. Si $g_i$ (resp. $h_i$) est le morphisme d'injection $X_i\to X$
(resp. $Y_i\to Y$), on en conclut que $f\circ g_i=h_i\circ f_i$, où $f_i$
est un morphisme dominant $X_i\to Y_i$
(\hyperref[I.5.2.2-fr]{I, 5.2.2})~; on est alors dans les conditions
d'application de (\hyperref[II.5.4.5-fr]{5.4.5}), et pour que $f$ soit
propre, il faut et il suffit que les $f_i$ le soient.
\end{env}

\begin{corollary}[5.4.8]
\label{II.5.4.8-fr}
Soient $X$, $Y$ des $S$-préschémas séparés de type fini sur $S$, et soit
$f:X\to Y$ un $S$-morphisme. Pour que $f$ soit propre, il faut et il suffit
que, pour tout $S$-préschéma $S'$, le morphisme
$f\times_S1_{S'}:X\times_SS'\to Y\times_SS'$ soit fermé.
\end{corollary}

Notons d'abord que si $g:X\to S$ et $h:Y\to S$ sont les morphismes
structuraux, on a par définition $g=h\circ f$, donc $f$ est séparé et de
type fini (\hyperref[I.5.5.1-fr]{I, 5.5.1} et
\hyperref[I.6.3.4-fr]{6.3.4}). Si $f$ est propre, il en est de même de
$f\times_S1_{S'}$ (\hyperref[II.5.4.2-fr]{5.4.2})~; a fortiori,
$f\times_S1_{S'}$ est fermé. Inversement, supposons vérifiée la condition de
l'énoncé et soit $Y'$ un $Y$-préschéma~; $Y'$ peut aussi être considéré
comme un $S$-préschéma, et comme $Y\to S$ est séparé, $X\times_YY'$
s'identifie à un sous-préschéma fermé de $X\times_SY'$
(\hyperref[I.5.4.2-fr]{I, 5.4.2}). Dans le diagramme commutatif
\[
  \xymatrix{
    X\times_YY'\ar[r]^{f\times1_{Y'}}\ar[d] &
    Y\times_YY'=Y'\ar[d]\\
    X\times_SY'\ar[r]_{f\times_S1_{Y'}} & Y\times_SY'
  }
\]
les flèches verticales sont des immersions fermées~; il en résulte aussitôt
que si $f\times_S1_{Y'}$ est un morphisme fermé, il en est de même de
$f\times_Y1_{Y'}$.

\begin{remark}[5.4.9]
\label{II.5.4.9-fr}
Nous dirons qu'un morphisme $f:X\to Y$ est \emph{universellement fermé} s'il
satisfait à la condition b) de la définition
(\hyperref[II.5.4.1-fr]{5.4.1}). Le lecteur observera que

